\section{Introduction}\label{sec:intro}

The Federal Reserve departed from its ``Treasuries-only'' framework in 2008 and 2020, purchasing trillions in agency MBS that suppressed 30-year mortgage rates to lows of roughly 2.6\% \citep{hancock2011} and triggered a refinance boom. In 2022 it reversed course. Quantitative Tightening (QT) was a passive roll-off intended to drain up to \$35 billion per month from its SOMA portfolio, under a cap phased in at \$17.5 billion and stepped up three months later. (June 2024 and April 2025 runoff slowings cut only the Treasury cap, and above-cap reinvestment never triggered in this window. So the constant post-phase-in \$35 billion MBS schedule is policy-committed throughout.) Unlike standard Treasuries, U.S. mortgages carry asymmetric, state-dependent payoff profiles. Tightening pushed the 30-year rate past 7.0\%, and under the par-payoff rule that imposed a substantial penalty on relocation: the lock-in effect. In practice, historic housing-inventory shortages and depressed consumer sentiment compounded it. Mortgages did not prepay.

Cumulative tightening-cycle runoff fell \$764.7 billion short of the phased caps \citep[on the balance-sheet setting, see][]{lopezsalido2023,acharya2023} (Section~\ref{sec:method-benchmark}). New York Fed staff projected, from the month the \$35 billion ceiling took effect, that the caps would sit above achievable runoff \citep{nyfed2022}: the ceiling was set roughly 1.7 to 1.9 times the Federal Reserve's own contemporaneous projection of achievable runoff. The observation is theirs; I add its quantification and ex-ante construction. Measured against that projection rather than against the never-binding caps, the shortfall is \$87.8 billion, 11.5\% of the cap-relative benchmark (Section~\ref{sec:method-benchmark}). The New York Fed's May 2022 staff baseline already anticipated the large majority of the realized cap-shortfall: 75.6--88.5\%, depending on a disclosed intra-2022 allocation choice.

That the lock-in effect slowed the runoff of the Federal Reserve's mortgage portfolio is not a new claim; its own staff and officials have said as much (\citealp{na2024}; \citealp{perli2024}; \citealp{hammack2025}). This paper's contribution is not that mechanism but its decomposition: an accounting exercise within a fully disclosed model, not a causal identification. Switching the lock-in elasticity off partitions the realized shortfall between an elastic channel and a rate-inelastic baseline: scheduled amortization plus baseline turnover held at its observed level. Its floor is read off realized, partly behavioral turnover, so the baseline is not a non-behavioral object. Most of the shortfall sits in that baseline: with the elasticity off, the loan-level model still recovers 85.7\% of the \$764.7 billion shortfall-against-cap benchmark at the calibration this paper headlines and 88.7\% in-sample, on the benchmark-consistent accounting basis (Section~\ref{sec:hazard-interp}). The Federal Reserve's own ex-ante projection implies the same composition, though not independently: it may already have embedded part of the lock-in channel. That 75.6--88.5\% range is wide enough that its overlap with the rate-inelastic null's recovery carries no precision.

What lock-in itself adds is secondary, and a range rather than a number. The elasticity's identified footprint is a bounded lock-in marginal, not a pinned magnitude and not a sign, and its causal content is inherited from the external estimates calibrating it, not established here (Section~\ref{sec:identification}). The floor that dominates the level is calibrated in-window, on the same locked-in cohorts the marginal decomposes. I therefore re-measure it outside that window, on pre-episode 2017--2019 discount cohorts, and re-run both legs at the result. The lock-in marginal runs $+4.3$ to $+6.8$ points across the off-window floor's defensible range. It runs $+2.3$ to $+9.1$ points once that floor's own sampling error is propagated with a wild-cluster small-sample correction and the pre-committed coverage rule selects the qualifying construction. That is the binding interval (Section~\ref{sec:robustness-floor}). At the mid-grid anchor inside that interval the value is $+5.6$ points, or $+\$42.6$ billion. The design delivers the range, not that value. In flow terms the lock-in marginal is roughly \$1 billion per month against a \$35 billion monthly redemption cap. Measured corrections to the floor and the accounting basis run in both directions: calendar-standardizing the floor read moves the value to $+4.7$ points; matching that read's housing-activity level to the window's moves it to $+10.7$. The in-window production calibration returns $+9.2$ points, reported as the in-sample calibration point, not the headline; its sampling intervals are in Table~\ref{tab:headline}, and the operative uncertainty remains calibration (Table~\ref{tab:uncertainty}). Neither figure is the entire shortfall.

A pre-committed sweep of the baseline turnover floor and the elasticity band bounds the lock-in marginal between $+2.1$ and $+13.2$ points, and it stays positive at every point of the calibration box and at every off-window anchor tested. That positivity checks the specification; it does not identify the mechanism. The grid, parity gates, and interpretive thresholds were committed to the repository before any run executed. I reserve ``pre-registered'' for third-party registries. That pre-commitment is a disclosure device, not a proof, and should be read net of the replication package's verdict appendix, which tabulates every such rule together with how its outcome was adjudicated after the run \citep{olken2015}.

I quantify this shortfall's mechanics with two structurally distinct estimators of the aggregate level (Sections~\ref{sec:abm} and~\ref{sec:hazard}). The first is a 10,000-household Agent-Based Model with dynamic macroeconomic frictions. The second is a loan-level hazard framework that abandons the household utility function altogether, introduced after the ABM's residual proved persistent and unresponsive to added behavioral realism. The ABM recovers 13.6\% on the fifty-seed mean (10--17\% across seed noise; 11.9\% on the frozen production draw), the hazard framework 91.3\% at the headline off-window floor and 97.9\% at the in-sample calibration point (remaining variants: Table~\ref{tab:headline}). The decomposition rests on the hazard framework alone: the rate-inelastic/rate-elastic partition is a hazard-framework object, because Path A is excluded from every headline range by pre-committed specification and is carried as appendix material (Section~\ref{sec:patha}), and the ABM's null admits two readings that disagree in sign. A cross-design variant of the ABM recalibrated on real loan covariates recovers 59.3\% on the frozen seed, 60.2\% across fifty seeds, and 76.3\% at the book's composition (Section~\ref{sec:robustness-crossdesign}), so the contrast is not clean evidence about modeling paradigm alone. The hazard framework is the instrument.

The par-payoff obligation sits inside a two-sided exchange. The pool homogeneity behind To-Be-Announced (TBA) market liquidity \citep[the market's institutional design is reviewed in][]{fuster2023} is purchased with household mobility, because the standardized contracts TBA delivery presupposes carry that obligation: to relocate, a homeowner must extinguish the existing mortgage at full face value. The obligation is a property of the U.S. fixed-rate note, not of TBA eligibility, and is enforced by the due-on-sale clause (made broadly enforceable by the Garn--St.~Germain Act of 1982). The FHA/VA loans backing Ginnie Mae pools are assumable, a statutory carve-out whose exercise is underwritten (HUD Handbook 4000.1; 38 U.S.C.~\S~3714). Those pools are 20.4\% of SOMA book face on the paper's CUSIP-level cohort parse, and this design nowhere measures their realized take-up, so that 20.4\% share bounds the carve-out from above. Observed speeds bound it a second time in the same direction: an assumed loan does not prepay, so materially high take-up would show up as \emph{slower} Ginnie voluntary speeds, and over the window Ginnie's voluntary component instead runs $1.35$ points per year \emph{above} the conventional one (Section~\ref{sec:pathb}). Both bounds are ceilings, not sizes. The decomposition disciplines the exchange: household mobility is constrained, as the household-level lock-in literature documents, while the cash-flow shortfall is governed mostly by loan-level and institutional mechanics, so the trade-off's binding cost is denominated in mobility, not institutional cash flow (Section~\ref{sec:conclusion}). That ranking rests on external evidence, not a within-paper comparison: I measure the institutional leg and import the mobility leg at an elasticity I never re-measure, so the two are nowhere commensurable. The one household-side magnitude I carry is \citepos{batzer2024} roughly \$2.4 trillion of below-market financing value forfeited on relocation, computed on universal relocation and therefore an upper bound. The ranking's direction is what that supports; its ratio is not.

This paper's affirmative contributions are three. First, the cap-design arithmetic: both components of the rate-inelastic path are computable in advance of the cycle, as a band (Sections~\ref{sec:identification} and~\ref{sec:discussion-capdesign}). Second, an expectations-based complement that measures the shortfall against the Federal Reserve's own ex-ante projection, not the never-binding caps (Section~\ref{sec:method-benchmark}). Third, a decomposition of the QT mortgage-runoff shortfall into rate-inelastic and rate-elastic components, the elastic part delivered as a bounded lock-in marginal (Sections~\ref{sec:abm}, \ref{sec:hazard} and~\ref{sec:identification}). Two further design elements are not contributions. The paired cross-design test is a validation device under interpretive thresholds fixed ex ante, and it ran against this paper's own ABM reading: on all fifty seeds the real-covariate leg crossed the pre-fixed threshold at which that reading is undercut (Section~\ref{sec:robustness-crossdesign}). The Danish market-value counterfactual and its incidence bracket are appendix material: a robustness result on a counterfactual, not a measurement of the U.S. book (Section~\ref{sec:abm-danish}; Appendix~\ref{app:params}).

Table~\ref{tab:headline} collects the core results and their operative uncertainties in one place; every value repeats a committed, gated figure.

{\footnotesize
\begin{longtable}{@{}p{4.6cm}p{4.6cm}p{5.3cm}@{}}
\caption{Core results and operative uncertainties. Each row repeats committed values derived in the referenced exhibit; recovery percentages are quoted on the benchmark-consistent shared basis. The uncertainty hierarchy is established in Table~\ref{tab:uncertainty}: simulation-seed noise is negligible at reporting precision and sampling is quoted under the stratum-cluster scheme; the calibration box is the operative uncertainty --- except the headline marginal, whose binding layer is the floor read's own propagated sampling error (Table~\ref{tab:uncertainty}). The replication package ledger's crosswalk table maps each headline quantity to its accounting basis and source run. The Path B recovery is estimated on the Freddie 2017--2021 conventional universe; Section~\ref{sec:method-benchmark} derives that universe's coverage of SOMA book face on the book's own agency $\times$ vintage joint cells, and states the external-speed overlay that bounds each exclusion (Section~\ref{sec:pathb}).}
\label{tab:headline}\\
\toprule
Object & Headline value & Uncertainty / bounds \\
\midrule
\endfirsthead
\toprule
Object & Headline value & Uncertainty / bounds \\
\midrule
\endhead
Cap-relative benchmark, June 2022--November 2025 (Section~\ref{sec:method-benchmark}) & \$764.7 billion & -- \\
Expectations-based complement (Section~\ref{sec:method-benchmark}) & \$87.8 billion, 11.5\% of the cap-relative benchmark & Fed's ex-ante projection implied 75.6--88.5\% of the cap-shortfall, by intra-2022 allocation \\
Rate-inelastic baseline, $\beta_1 = 0$ (Section~\ref{sec:identification}; Table~\ref{tab:bases}) & 85.7\% of benchmark at the headline off-window floor; 88.7\% in-sample, both under the production floor form (35.6\% under the additive form) & 83.9--86.9\% across the off-window floor range under that form; moves one-for-one with the floor level; floor-read wild-cluster $[80.6, 88.4]$\% (run \texttt{null\_floor\_interval}) \\
Lock-in marginal, off-window floor (Sections~\ref{sec:identification} and~\ref{sec:robustness-floor}; Table~\ref{tab:oosfloor}) & $+2.3$ to $+9.1$ points, basis-invariant --- the floor read's own sampling error propagated at the central elasticity (the 6.5\% mid-band value), not an interval on the elasticity, whose dimension Table~\ref{tab:lowband} sweeps; $+5.6$ points $= +\$42.6$ billion at the mid-grid anchor inside it, an anchor convention rather than a central tendency; the range carries the identified content (Section~\ref{sec:identification}) & floor-read wild-cluster bootstrap-$t$ inversion $[+2.3, +9.1]$, the binding layer (31 clusters; percentile read $[+3.0, +8.0]$, under-covering; the corrected ladder and its degrees of freedom are Table~\ref{tab:ladder}) \\
Lock-in marginal, convention envelope (Table~\ref{tab:assembly}) & $+0.9$ to $+15.6$ points across named convention variants (seasoning ramp 75--150 PSA; moving-share bracket; housing-activity-scaled baseline; floor form) & convention ranges without coverage properties; the calibration box spans $+2.1$ to $+13.2$ \\
\addlinespace
Lock-in marginal, in-sample calibration point (Section~\ref{sec:identification}; Table~\ref{tab:uncertainty}) & $+9.2$ points $= +\$70.3$ billion, basis-invariant & sampling: within-stratum $[+9.17, +9.23]$, stratum-cluster $[+8.27, +10.19]$ (the operative scheme); calibration box $+2.1$ to $+13.2$; Fannie replication $+8.68$ (Section~\ref{sec:pathb-fannie}) \\
Path B central level (Table~\ref{tab:estimators}) & 91.3\% at the headline off-window floor; 97.9\% in-sample (100.0\% composed as committed, 98.1\% on a single coupon convention) & 88.2--93.7\% across the off-window floor range; covariate priors 92.2--99.2\% (in-sample); Ginnie-speed overlay at the headline floor 84.0\% central against a 79.5\% null, and 89.3\% at the in-sample point (raw; of the \$66.3B raw adjustment the Ginnie-specific piece is $+\$38.5$B, within the \$20--47B bound) \\
Production ABM (Table~\ref{tab:estimators}) & 13.6\% seed mean & 10--17\% seed band; frozen draw 11.9\%; the operative uncertainty is calibration, not seed noise --- established on the real-covariate cross-design, where one calibration choice moves recovery across 20.9--59.3\% (Section~\ref{sec:robustness-crossdesign}); that span belongs to two calibration choices on that design, not to this row's 13.6\%, and does not contain it \\
Cross-design ABM, real covariates (Table~\ref{tab:crossdesign}) & 59.3\% recalibrated; fifty-seed mean 60.2\%, 95\% band 55.1--66.0\% (76.3\% at book composition; effective sample 1{,}332 of 10{,}000, re-draw 70.6\%) & 20.9\% under the frozen scale; fifty-seed mean 22.0\% (12.6\% at book composition) \\
\addlinespace
Portfolio duration extension (Table~\ref{tab:wal}; Appendix~\ref{sec:robustness-wal}) & empirical-path WAL 9.4 years at the window's open and 8.5 at its close --- a 6.0-year extension against the no-shock counterfactual's 2021 refinancing-boom speeds (3.4 years) & the comparator sets the result, and this row is quoted against both. Read instead against the baseline turnover floor's off-window read, the same path runs 9.4 years against 9.5 at the open and 8.5 against 8.6 at the close, so the six years measure refinancing-versus-turnover rather than lock-in-versus-normalcy. Approximate single-pool-per-cell calculation; option-adjusted duration and convexity are declined as out of scope \\
\bottomrule
\end{longtable}
}

\paragraph{Definitions used throughout.} Four accounting bases attach to every recovery percentage (Section~\ref{sec:method-benchmark}): \emph{standalone-scorer} (no adjustment), \emph{benchmark-consistent shared} (nets a common \$69.6 billion curtailment flow identically from every U.S. leg; the main-text convention), \emph{full-book} (reweights to the SOMA book's coupon mix), and \emph{composed} (both). Two floor calibrations attach to every marginal and level: the \emph{in-sample calibration point} (the 4\% baseline turnover floor, in-window on the locked-in 2023--24 cohorts; returns $+9.2$ points) and the \emph{off-window calibration} (the 4.70--5.33\% floor, on pre-episode 2018 discount cohorts; returns $+5.6$ points and is the headline). The floor is a measured object: realized-turnover reads put it at 3.8--3.9\% in-window (flat across seasoning cuts); 4.70--5.33\% on the off-window 2018 leg, all of it on cohort-months aged twelve to twenty-four, so that leg's mature-seasoning test is not computable rather than passed; 5.51\% under an imputation-heavy age standardization; and 5.52\% on an out-of-agency Fannie Mae read of the pooled off-window period---the last two sit above the clean band, hence the band's soft lower edge (Section~\ref{sec:robustness-floor}). The ``baseline turnover floor'' (sometimes called an involuntary-turnover floor) names that object without decomposing it: the reads measure total deep-discount-cohort turnover---discretionary life-cycle moves and cash-out refinancings mixed with strictly involuntary events---and the strictly-involuntary share, plausibly well under half (Section~\ref{sec:pathb}), is priced directly by the mixture curve of Section~\ref{sec:robustness-floor}. ``Trapped liquidity'' and ``shortfall'' denote the net shortfall of realized roll-off against the phased redemption caps (\$764.7 billion, June 2022--November 2025). The \emph{lock-in marginal}, the paper's identified object, is the central simulation minus the \emph{rate-inelastic baseline}, its $\beta_1 = 0$ null: scheduled amortization from the book's own coupon, age, and term composition plus baseline turnover held at its observed level, the marginal being the imported elasticity's increment above it. The partition separates baseline from increment, not mechanics from behavior. The marginal's operative uncertainty is the calibration-and-form envelope of Table~\ref{tab:uncertainty} rather than any sampling interval, and the envelope is a grid, not a confidence region---a load-bearing distinction: the elasticity band it sweeps is \citeauthor{liebersohn2024}'s \emph{specification} range with my adopted midpoint at its centre, not a published point estimate with a published standard error, and so the envelope has no coverage property. Finally, \emph{no outcome-holdout months exist anywhere in this paper}, without exception: every recovery percentage is in-sample or calibrated in the time dimension, a description of fit rather than a prediction (Section~\ref{sec:limitations}).
